% !TeX root = ../main.tex
\section{Computational details of Bayesian-PMALA}
\label{app:pmala}

\subsection{Full conditional distributions}
\label{app:pmala:conditionals}

Write $e_i = y_i - \theta_i$. For $z_i$, expanding the square in \eqref{eq:posterior} and dropping the cross term $-\psi_\tau e_i/(\omega_\tau^2\sigma)$, which does not depend on $z_i$, gives
\[
p(z_i \mid \theta, \sigma, \by) \propto z_i^{-1/2} \exp\!\left\{-\frac{1}{2}\left(\frac{\chi_i}{z_i} + \tilde\psi\, z_i\right)\right\},
\qquad \chi_i = \frac{e_i^2}{\omega_\tau^2 \sigma}, \quad \tilde\psi = \frac{1}{\sigma}\left(\frac{\psi_\tau^2}{\omega_\tau^2} + 2\right),
\]
a generalized inverse Gaussian distribution $\mathrm{GIG}(1/2, \chi_i, \tilde\psi)$, with $\tilde\psi = 1/\{2\tau(1-\tau)\sigma\}$. For $e_i \ne 0$ its reciprocal is inverse Gaussian with mean $(\tilde\psi/\chi_i)^{1/2} = 1/\{\tau(1-\tau)|e_i|\}$ and shape $\tilde\psi$, which is \eqref{eq:zupdate}; for $e_i = 0$, $z_i$ is Gamma distributed with shape $1/2$ and rate $\tilde\psi/2$, a case that safeguard~(i) of Appendix~\ref{app:pmala:safeguards} excludes. The terms involving $\lambda$ are $\lambda^{a_0+m-1}\exp\{-\lambda(b_0 + \norm{D\theta}_1)\}$, which gives \eqref{eq:lambdaupdate}. The factors $\sigma^{-n/2}$ and $\sigma^{-n}$ from the Gaussian and exponential densities and the inverse-Gamma prior give \eqref{eq:sigmaupdate}.

\subsection{The proximal map of the individual sampler}
\label{app:pmala:admm}

Write the metric \eqref{eq:metric} as $M = \mathrm{diag}(d) + cD^\top D$, with $d = w$ and $c = \min(\kappa\lambda^2/2, 10^8)$, the cap being safeguard~(ii) of Appendix~\ref{app:pmala:safeguards}. Problem \eqref{eq:prox} is the minimization over $(\theta, v)$ of $\frac12\sum_i w_i(\theta_i - r_i)^2 + \lambda\norm{v}_1 + \norm{\theta - x}_M^2/(2\gamma)$ subject to $v = D\theta$. We apply scaled ADMM with penalty $\varrho > 0$ and scaled dual variable $\nu \in \R^m$. Starting from $v^0 = Dx$ and $\nu^0 = 0$, iteration $t$ is
\begin{align*}
\theta^{t+1} &= x + A^{-1}\left\{ W(r - x) + \varrho D^\top (v^t - \nu^t - Dx) \right\}, \qquad A = W + \gamma^{-1} M + \varrho D^\top D, \\
v^{t+1} &= \mathcal{S}_{\lambda/\varrho}\!\left(D\theta^{t+1} + \nu^t\right), \qquad
\nu^{t+1} = \nu^t + D\theta^{t+1} - v^{t+1},
\end{align*}
where $\mathcal{S}_a(b) = \mathrm{sign}(b)\max(|b| - a, 0)$ acts elementwise. The $\theta$-update is the usual one, $A\theta^{t+1} = Wr + \gamma^{-1}Mx + \varrho D^\top(v^t - \nu^t)$, solved for the increment $\theta^{t+1} - x$: its right-hand side scales with the move $T(x) - x$ rather than with $\gamma^{-1}Mx$, so its rounding error does not grow like $1/\gamma$ in the acceptance ratio \eqref{eq:pmala_accept}. We use $\varrho = 0.1\,(\bar w + c)/\gamma$ with $\bar w = n^{-1}\sum_i w_i$, about one tenth of the curvature of the metric term, and return $\theta^{150}$, with no data-dependent stopping rule.

The matrices $A$ and $M$ are symmetric positive definite and banded with half-bandwidth $k+1$. Within one move $A$ does not depend on $x$, so one banded Cholesky factor serves the forward map $T(\theta)$ and the reverse map $T(\theta^\star)$. The metric is factorized once per move as $M = LL^\top$ for the noise in \eqref{eq:pmala_proposal}, and the norms in \eqref{eq:pmala_accept} are evaluated as $\norm{b}_M^2 = \sum_i d_i b_i^2 + c\norm{Db}_2^2$. A factorization costs $O(n(k+1)^2)$ operations and an ADMM iteration $O(n(k+1))$.

The correction \eqref{eq:pmala_accept} requires only that $T$ be a deterministic function of its input, given $(z, \lambda, \sigma, \delta)$ (Section~\ref{subsec:pmala}). The cold start from $(Dx, 0)$ and the fixed iteration rule, used for the forward and the reverse map, guarantee this, whereas a warm start from the previous sweep would not. For the same reason the noise is added after the proximal map: added before it, it would give the image of a Gaussian under a nonlinear map, whose density is not available.

\subsection{The proximal map of the joint sampler}
\label{app:pmala:boxsolver}

Fix a level and suppress its index. The box $B = [\ell, u]$ of \eqref{eq:joint_box} may have infinite components (at $j = 1$ and $j = J$), and $\ell_i = u_i$ is allowed where two neighboring curves touch; such coordinates are fixed at their bound and eliminated. We write $T_B$ for the map \eqref{eq:joint_prox}, the proximal map of $U + \iota_B$ in the $M$-norm. The smooth part of its objective has Hessian $A_0 = W + M/\gamma$, so on the free coordinates the problem is strongly convex with modulus at least $\varpi = \min_{i:\,\ell_i < u_i}(w_i + d_i/\gamma)$. We write $\widehat T_B$ for the map actually computed, the $T_j$ of \eqref{eq:joint_proposal}--\eqref{eq:joint_accept}.

\paragraph{Stage 1: the unbounded map.} The solver first computes $\tilde\theta = T(x)$ by the ADMM of Appendix~\ref{app:pmala:admm}, with the same penalty, cold start and budget, and returns this iterate, the output of a fixed number of ADMM iterations, whenever it is feasible, $\ell \le \tilde\theta \le u$. The rule mirrors the fact that an exact unbounded minimizer lying in $B$ is also the constrained minimizer. Hence $\widehat T_B$ coincides bitwise with the individual map for a free box ($J = 1$) and whenever the computed $\tilde\theta$ lies in $B$.

\paragraph{Stage 2: an interior-point method.} Otherwise $T_B(x)$ is computed from the epigraph quadratic program
\begin{equation}
\begin{aligned}
\min_{\theta \in \R^n,\ \upsilon \in \R^m}\ & \tfrac12\theta^\top A_0\theta + \beta^\top\theta + \lambda\mathbf 1^\top\upsilon \\
\text{subject to}\ & D\theta - \upsilon \le 0, \quad -D\theta - \upsilon \le 0, \quad \theta_i \le u_i\ (u_i < \infty), \quad \ell_i - \theta_i \le 0\ (\ell_i > -\infty),
\end{aligned}
\label{eq:box_qp}
\end{equation}
with $\beta = -(Wr + Mx/\gamma)$, which is equivalent to \eqref{eq:joint_prox} because $\upsilon = |D\theta|$ at the optimum. We solve \eqref{eq:box_qp} by Mehrotra's predictor--corrector primal--dual method \citep{mehrotra1992implementation}, with slacks $s_1, \ldots, s_4$ and multipliers $\mu_1, \ldots, \mu_4$ for the four constraint blocks. Eliminating the slack and multiplier steps and, through the $\upsilon$-row of the Newton equations, the step in $\upsilon$, leaves one linear system in $\Delta\theta$ with matrix $A_0 + D^\top \Lambda_\upsilon D + \Lambda_3 + \Lambda_4$, where $\Lambda_l = \mathrm{diag}(\mu_l/s_l)$, $l = 1, \ldots, 4$, and $\Lambda_\upsilon = 4\Lambda_1\Lambda_2(\Lambda_1 + \Lambda_2)^{-1}$, which is symmetric positive definite and banded with half-bandwidth $k+1$. In each iteration the predictor and the corrector each factorize it by banded Cholesky, in $O(n(k+1)^2)$ operations, and solve with it, with centering parameter $\min\{1, (\bar\mu_{\mathrm{aff}}/\bar\mu)^3\}$, where $\bar\mu = s^\top\mu/N$ over the $N$ inequality constraints and $\bar\mu_{\mathrm{aff}}$ is its value after the predictor step. Separate primal and dual step lengths in $(0, 1]$ keep $0.99$ of the distance to the boundary.

\paragraph{Initialization and stopping.} The method starts from $\theta^0 = \min\{\max(\tilde\theta, \ell), u\}$, $\upsilon^0 = |D\theta^0| + s_{\min}$, slacks at least $s_{\min} = 0.01\,s_\theta$, $s_\theta = \max(1, \norm{\theta^0}_\infty)$, $\mu_1 = \mu_2 = \lambda/2 + 10^{-8}$ and box multipliers $1$. With $\varepsilon = 10^{-8}$, the thresholds are $\varepsilon(1 + s_\theta)$ for the primal residual, $\varepsilon\varpi$ for the dual residual and $\varepsilon^2\varpi/2$ for $s^\top\mu$. With $\phi$ the largest ratio of a residual to its threshold, the iterate with the smallest $\phi$ is retained, clipped to $B$ and returned when $\phi \le 1$, on stall ($10$ iterations after the $15$th without a $10\%$ reduction of the best $\phi$), on a non-finite value, or after $100$ iterations. In the sampler $\varepsilon\varpi$ is at the rounding level of $A_0\theta + \beta$, so the loop usually ends on stall. These are stopping criteria, not error bounds; the accuracy of the output is checked in Appendix~\ref{app:pmala:tests}.

\paragraph{Determinism.} Every rule above depends only on $(x, w, r, \lambda, d, c, \gamma, \ell, u)$ and on the iterates, so $\widehat T_B$ is a deterministic function of $x$ and the correction \eqref{eq:joint_accept} is exact for the computed map. The forward and reverse maps of a move use the same box, metric, step size and rule.

\subsection{Interweaving step for the penalty}
\label{app:pmala:asis}

The conjugate draw \eqref{eq:lambdaupdate} is followed by a draw in the ancillary parameterization \citep{yu2011center}. Let $Q \in \R^{n\times(k+1)}$ be an orthonormal basis of the null space of $D$, and write $\theta = p + v$ with $p = QQ^\top\theta$. Under \eqref{eq:prior_theta} the standardized differences $\lambda D\theta = \lambda Dv$ have independent standard Laplace entries and do not depend on $\lambda$. Holding them and $(p, z, \sigma)$ fixed, a new value $\lambda'$ implies the trend $\theta(\lambda') = p + (\lambda/\lambda')v$. The factor $\lambda'^m$ in \eqref{eq:prior_theta} cancels the Jacobian of $v \mapsto \lambda' Dv$, so $\eta = \log\lambda'$ has log density
\begin{equation}
\log f(\eta) = a_0\eta - b_0 e^\eta - \frac12\sum_{i=1}^n w_i\left\{\theta_i(e^\eta) - r_i\right\}^2 ,
\label{eq:asis_density}
\end{equation}
where $a_0\eta$ includes the factor $e^\eta$ from the change of variables. We draw $\eta$ by slice sampling \citep{neal2003slice} with initial width $1$, at most $30$ stepping-out steps divided at random between the two ends, and shrinkage, leaving the state unchanged if the interval shrinks below width $10^{-10}$, and set $\theta \leftarrow \theta(\lambda')$. The conjugate and the ancillary draws each leave \eqref{eq:posterior} invariant, and so does their composition.

\paragraph{Constrained version for the joint sampler.} At level $j$, with $p$ and $v$ as in \eqref{eq:joint_asis_curve} and the index dropped, the rescaled trend $p + \varphi v$, $\varphi = \lambda/\lambda' > 0$, lies in $B_j$ if and only if $\ell_i \le p_i + \varphi v_i \le u_i$ for every $i$. For $v_i \ne 0$ this is an interval in $\varphi$ with endpoints $(\ell_i - p_i)/v_i$ and $(u_i - p_i)/v_i$, and for $v_i = 0$ it requires $p_i \in [\ell_i, u_i]$. The intersection of these intervals with $(0, \infty)$ is $[\varphi_-, \varphi_+]$, with endpoints $0$ or $\infty$ allowed, and
\begin{equation}
\log\lambda - \log\varphi_+\ \le\ \eta\ \le\ \log\lambda - \log\varphi_- .
\label{eq:joint_asis_interval}
\end{equation}
The slice update targets \eqref{eq:asis_density} restricted to \eqref{eq:joint_asis_interval}, with the weights and pseudo-responses of the newly drawn $z_j$ and the current $\sigma_j$. The reconstructed trend is also checked against the box, so that rounding cannot produce a crossing, and is never clipped. If the interval is empty or degenerate, or the current density is not finite, the state is left unchanged.

\subsection{Initialization, adaptation and safeguards}
\label{app:pmala:init}
\label{app:pmala:adapt}
\label{app:pmala:safeguards}

\paragraph{Initialization and warm-up.} In the algorithm of Section~\ref{sec:methodology}, every chain of the individual sampler starts at $\theta^{(0)} = \by$, with $\sigma^{(0)} = 1$, $z_i^{(0)} = \sigma^{(0)}$, $\lambda^{(0)} = m$ and $\delta^{(0)} = 0.2$. These data-free values avoid the conditionals given $\theta = \by$, which place $\sigma$ and $z$ near zero, a configuration from which a Gibbs chain moves away only slowly. For $J > 1$ the joint sampler starts the trends at $\theta_j^{(0)} = \by + s_y\{\Phi^{-1}(\tau_j) - \mathrm{median}_{h}\,\Phi^{-1}(\tau_h)\}\mathbf 1$, with $s_y = \max\{\mathrm{sd}(\by), 0.1\}$ and $\Phi$ the standard normal distribution function, so that the levels start strictly ordered; the other parameters start as in the individual sampler.

From $\theta = \by$ an exact Metropolis--Hastings chain barely moves: proposals are close to the output of a proximal map, which is smooth in $D\theta$, so for the rough $\by$ the reverse term $\norm{\by - T(\theta^\star)}_M^2/(2\delta)$ in \eqref{eq:pmala_accept} can make the acceptance probability of a proposal negligible. The first $H = 200$ sweeps therefore replace the P-MALA move by the proximal-point step $\theta \leftarrow \arg\min_{\vartheta}\{U(\vartheta) + \norm{\vartheta - \theta}_W^2/(2\gamma_{\mathrm{w}})\}$, $\gamma_{\mathrm{w}} = 10^4$, approximated by $300$ iterations of the ADMM of Appendix~\ref{app:pmala:admm} with $M = W$, $\gamma = \gamma_{\mathrm{w}}$ and $\varrho = \mathrm{median}(w_i)$. In the joint sampler the step is restricted to the current box, that is \eqref{eq:joint_prox} with $M_j = W_j$ and $\gamma_{\mathrm{w}}$, and computed as in Appendix~\ref{app:pmala:boxsolver} with the same budget, which keeps the levels ordered. The blocks $z$, $\lambda$ and $\sigma$ are drawn as usual during the warm-up; the interweaving step and the step-size adaptation start after it. The warm-up is part of the burn-in.

\paragraph{Step-size adaptation.} After the warm-up the step size is updated at every sweep by $\log\delta_{t+1} = \log\delta_t + t^{-0.6}\{\min(1, \alpha_t) - 0.5\}$, with $\alpha_t$ from \eqref{eq:pmala_accept}, and $\delta$ is confined to $[10^{-8}, 10^{2}]$. The adaptation continues after burn-in, because $\lambda$ can still drift slowly in the retained sweeps and a frozen step size would let the acceptance rate drift with it; the vanishing gain satisfies the diminishing-adaptation condition of \citet{roberts2007coupling}. In the joint sampler each level has its own $\delta_j$, and a proposal outside $B_j$ counts as $\alpha_t = 0$.

\paragraph{Safeguards.} (i) Each draw $z_i$ is restricted to $[10^{-8}, 10^{4}]$, and $|y_i - \theta_i|$ is floored at $10^{-10}$ in \eqref{eq:zupdate}, so that $M$ and the solver systems can be factorized. (ii) The coefficient $\kappa\lambda^2/2$ in \eqref{eq:metric} is capped at $10^{8}$. (iii) If a banded Cholesky factorization fails because of rounding, a ridge $\epsilon\max_i A_{ii}$ is added to the diagonal, starting from $\epsilon = 10^{-12}$ and multiplying $\epsilon$ by $100$ up to five times; the ridge is part of the solver, so the proximal map remains deterministic. (iv) A non-finite $\log\alpha$ is treated as a rejection, and a non-finite state stops the run with an error. The joint sampler also verifies after every update of a level that the trend lies in its box, and a violation stops the run.

The Metropolis--Hastings argument of Section~\ref{subsec:pmala} covers (ii) and (iii), which change only the proposal, and the box check, which never triggers when the kernels are implemented correctly. It does not cover (i): clipping the latent draws and flooring the residual alter the conditional of $z$ slightly, and clipping is not sampling from a normalized truncated distribution, so the sampler targets \eqref{eq:posterior} or \eqref{eq:joint_target} up to the clipping of the latent draws; the trend update itself remains exact.

\subsection{Settings of each experiment}
\label{app:pmala:settings}

The algorithm of Section~\ref{sec:methodology} uses $\kappa = 3$, $a_\sigma = b_\sigma = 0.1$, $\delta^{(0)} = 0.2$, $\lambda^{(0)} = m$, $\sigma^{(0)} = 1$, $H = 200$, $\gamma_{\mathrm{w}} = 10^4$, $150$ ADMM iterations, and gain $t^{-0.6}$ toward acceptance $0.5$ with $\delta \in [10^{-8}, 10^2]$. Table~\ref{tab:pmala_settings} lists, for each result, the settings recorded for its runs.

\begin{table}[htbp]
\centering
\footnotesize
\caption{Sampler settings of each experiment. Target: acceptance target. ADMM: form of the $\theta$-update (Appendix~\ref{app:pmala:admm}). ---: not recorded. $^{a}$Levels $0.96$--$0.99$. $^{b}$Levels $0.90$ and $0.93$, from a second run with the same sampler and kernel code, whose run configuration was not recorded. $^{c}$Prior $\lambda\sim\mathrm{Gamma}(n, 10^{-2})$, that is, $a_0 = n$; the SPod fits use two chains per level.}
\label{tab:pmala_settings}
\setlength{\tabcolsep}{2.5pt}
\begin{tabular}{@{}llllllll@{}}
\toprule
Result & Section & Version & $b_0$ (VS / S) & Sweeps/burn-in & Interweaving & Target & ADMM \\
\midrule
Interior & \ref{sec:sim:individual} & App.~\ref{app:pmala:interior} & $10^{-3}$ / $5\times10^{-4}$ & $30000/10000$ & spline-held & $0.5$ & increment \\
Extreme, individual$^{a}$ & \ref{sec:sim:extreme} & Sec.~\ref{sec:methodology} & $10^{-2}$ / $10^{-3}$ & $10000/3000$ & off & $0.55$ & increment \\
Extreme, individual$^{b}$ & \ref{sec:sim:extreme} & Sec.~\ref{sec:methodology} & --- & --- & --- & --- & increment \\
Extreme, joint & \ref{sec:sim:extreme} & Sec.~\ref{sec:methodology} & $10^{-2}$ / $10^{-3}$ & $30000/10000$ & on & $0.5$ & increment \\
Air simulation & \ref{sec:air_simulation} & App.~\ref{app:pmala:interior} & $10^{-2}$\,$^{c}$ & $30000/10000$ & spline-held & $0.5$ & increment \\
SPod data & \ref{sec:app:spod} & App.~\ref{app:pmala:interior} & $10^{-2}$\,$^{c}$ & $30000/10000$ & spline-held & $0.5$ & increment \\
Growth data & \ref{sec:app:growth} & --- & $\sqrt{G}$ & $6000/1000$ & --- & --- & --- \\
Timings & \ref{sec:scalability} & --- & --- & --- & --- & --- & --- \\
\bottomrule
\end{tabular}
\end{table}

\subsection{Version of the sampler used for the interior simulations and the air-quality fits}
\label{app:pmala:interior}

The interior-quantile results of Section~\ref{sec:sim:individual} use a version of the individual sampler that differs from the algorithm above in the starting values, the treatment of $\lambda$ during the warm-up, the interweaving step and two safeguards. The air-quality fits of Sections~\ref{sec:air_simulation} and~\ref{sec:app:spod} use the same version, with the prior $\lambda\sim\mathrm{Gamma}(n, 10^{-2})$.

\paragraph{Start and warm-up.} With window length $\max(5,\,2\lfloor n/40\rfloor+1)$, which is $25$ at $n = 500$, let $\theta^{\mathrm{rm}}$ be the running median of $\by$, and let $\hat q_\tau$ be the empirical $\tau$-quantile of $\by - \theta^{\mathrm{rm}}$. The chain starts at $\theta^{(0)} = \theta^{\mathrm{rm}} + \hat q_\tau\mathbf 1$, with $\sigma^{(0)} = n^{-1}\sum_i\rho_\tau(y_i - \theta^{(0)}_i)$, the maximum-likelihood ALD scale given $\theta^{(0)}$, with $z^{(0)}$ drawn from \eqref{eq:zupdate} given $(\theta^{(0)}, \sigma^{(0)})$, and with $\lambda^{(0)} = m$. In each warm-up sweep the proximal-point step is computed by the ADMM of Appendix~\ref{app:pmala:admm} if its relative residual falls below $10^{-6}$, and otherwise by the interior-point method of Appendix~\ref{app:pmala:boxsolver} with $\varepsilon = 10^{-8}$ and at most $200$ iterations; in the interior simulations the interior-point method was used at every warm-up sweep. In three replayed interior-simulation runs, in the $300$ air-quality simulation fits of Table~\ref{tab:air_n500} and in the six chains fitted to the SPod data, it was used at each of the $200$ warm-up sweeps and stopped before reaching its tolerance at every one of them, so the warm-up, whose purpose is to move the chain away from its start, applies an approximate proximal-point step and is discarded with the burn-in. The draws of $z$ and $\sigma$ are unchanged, but $\lambda$ is held at $\lambda^{(0)}$; its conjugate draw and the interweaving step start at sweep $H + 1$.

\paragraph{Interweaving step.} Let $P_q$ be the orthogonal projection onto the span of $q$ cubic B-splines with equally spaced knots on the design points $i/n$, a space that contains the polynomials of degree at most $k$, and write $\theta = p + v$ with $p = P_q\theta$. A new value $\lambda'$ gives $\theta(\lambda') = p + (\lambda/\lambda')v$, and $\eta = \log\lambda'$ is drawn by the slice sampler of Appendix~\ref{app:pmala:asis} from
\[
(a_0 + q - k - 1)\,\eta - b_0 e^{\eta} - \norm{e^{\eta}Dp + \lambda Dv}_1 - \frac12\sum_{i=1}^n w_i\left\{\theta_i(e^{\eta}) - r_i\right\}^2 .
\]
Its exponent collects $a_0 - 1 + m$ from the prior, $1$ from the change of variables and $-(n - q)$ from the rescaling of $v$. For $q > k+1$, $Dp$ can be nonzero, so the step is not an ancillary draw but a generalized Gibbs move \citep{liu2000generalised}: the transformations $(\lambda, \theta) \mapsto (e^{t}\lambda,\, p + e^{-t}v)$, $t \in \R$, form a group, the density above is the conditional of $\eta = \log\lambda + t$ that their construction prescribes, and the step leaves \eqref{eq:posterior} invariant. For $q = k+1$ it is the step of Appendix~\ref{app:pmala:asis}. The dimension $q$ cycles through $10$, $20$ and $40$, one value per sweep, and the fixed cycle also leaves the posterior invariant.

\paragraph{Safeguards and diagnostics.} In safeguard~(i) the floor on $|y_i - \theta_i|$ is $10^{-100}$, and the inverse Gaussian draw in \eqref{eq:zupdate} is computed in a form that avoids cancellation and has the same distribution. During the burn-in only, a draw of $\sigma$ larger than $10\sigma^{(0)}$ is replaced by $\sigma^{(0)}$. Over the $3000$ interior runs this reset was never triggered, no run stopped, and the post-burn-in acceptance rate was between $0.486$ and $0.520$. The median wall time per run was $153$ seconds on one core of a Linux compute cluster, and the median effective sample size of $\log\lambda$, from every $30$th retained draw, was $102$ for the VS signal and $132$ for the S signal. For the air-quality fits under the prior $\mathrm{Gamma}(n, 10^{-2})$ the trend mixes slowly; we report the posterior mean in the simulation and the posterior median for the SPod data as point estimates, and for the SPod data we pool two chains.

\subsection{Exactness checks}
\label{app:pmala:tests}

\paragraph{Individual sampler.} Tests check the identity \eqref{eq:ald_mixture}; the GIG conditional of $z_i$; the proximal map against the \texttt{glmgen} solver \citep{tibshirani2014adaptive} at $M = \mathbf{I}_n$, and its minimizing property in the metric \eqref{eq:metric}; the trend move against the exact Gaussian conditional at $\lambda = 0$ and against grid integration at $n = 4$; the shared-factorization move against separate solves; and the interweaving chain against grid integration. At $n = 7$ the step of Appendix~\ref{app:pmala:interior} with $q \in \{4, 5, 6\}$ agreed with an importance-sampling reference, which detected a deliberately wrong coefficient of $\eta$.

\paragraph{Joint sampler.} At $\lambda = c = 0$, $\widehat T_B$ matches the closed-form separable solution to $8.5\times10^{-12}$. Its objective was never above that of the solver Clarabel \citep{goulart2024clarabel} on $108$ small problems, and above it by at most $1.9\times10^{-9}$ relative on $72$ problems that reproduce the conditions of the sampler; on the $66$ that reached stage~2, a tighter tolerance moved the output by at most $4\times10^{-9}$ proposal-noise standard deviations, and an independent ADMM came within $0.033$ such standard deviations. With $J = 1$ and a common seed the joint sampler reproduces the individual sampler exactly on the same machine. At $n = 4$, with a box that cuts into the bulk of the conditional, the constrained trend move, alone and with the draw of $\lambda$ and the constrained interweaving step, agrees with integration on an $80^4$ midpoint grid (batch-means $z$-scores at most $1.9$; $E(\lambda) = 1.089 \pm 0.004$ against the exact $1.090$). A three-level run ($n = 80$, $\tau \in \{0.5, 0.9, 0.95\}$) has no crossing in any retained draw.
